% 似地行星（综述）
% license CCBYSA3
% type Wiki

本文根据 CC-BY-SA 协议转载翻译自维基百科\href{https://en.wikipedia.org/wiki/Earth_analog}{相关文章}。

\begin{figure}[ht]
\centering
\includegraphics[width=8cm]{./figures/849251199efbb53d.png}
\caption{地球与金星的演化路径。金星一直是一个典型范例，用来说明一颗外观上类似地球的行星，其演化结果可能会与地球产生何等巨大的差异。} \label{fig_Eaan_1}
\end{figure}
类地行星（Earth analog），也称为“地球双胞胎”或“第二地球”，是指其环境条件与地球相似的行星或卫星。“类地行星”（Earth-like planet）这一术语也常被使用，但该术语有时泛指任何类地行星。

这一可能性对天体生物学家和天文学家而言尤为重要，其基本逻辑在于：一颗行星与地球越相似，它就越有可能具备维持复杂地外生命的能力。因此，这一主题长期以来一直是科学、哲学、科幻文学以及大众文化中的重要猜想和讨论对象。空间殖民以及“太空与人类存续”理念的倡导者，也长期以来一直在寻找可供定居的类地行星。在遥远的未来，人类甚至可能通过行星改造（terraforming）来人工制造一颗类地行星。

在对系外行星进行科学搜索与研究之前，这一可能性主要通过哲学和科幻作品来探讨。哲学家曾提出，由于宇宙尺度极其巨大，某处必然存在一颗与地球几乎完全相同的行星。“平庸原理”认为，像地球这样的行星在宇宙中应当十分常见；而“稀有地球假说”则认为此类行星极其罕见。迄今为止已发现的成千上万个系外行星系统，在整体结构上与太阳系有着显著差异，这在一定程度上支持了稀有地球假说。

2013 年 11 月 4 日，天文学家基于开普勒空间望远镜的数据报告称，在银河系内围绕类太阳恒星和红矮星的宜居带中，可能存在多达 400 亿颗地球大小的行星 $^{[1][2]}$。从统计意义上看，距离地球最近的此类行星可能位于约 12 光年之内 $^{[1][2]}$。2020 年 9 月，天文学家在已确认的 4000 多颗系外行星中，依据天体物理参数以及地球生命的自然演化历史，识别出了 24 颗“超级宜居行星”（即在宜居性方面优于地球的行星）候选者 $^{[3]}$。
\begin{figure}[ht]
\centering
\includegraphics[width=8cm]{./figures/22fbf76c2e490700.png}
\caption{一颗假想的宜居系外行星及其三颗天然卫星的艺术示意图。} \label{fig_Eaan_2}
\end{figure}
自20世纪90年代以来的科学发现，极大地影响了天体生物学领域的研究范围、行星宜居性模型，以及对地外智能生命（SETI）的探索。
\subsection{历史}
\begin{figure}[ht]
\centering
\includegraphics[width=8cm]{./figures/4c512b83e6f4ad53.png}
\caption{珀西瓦尔·洛威尔将火星描绘成一颗干燥但类似地球、并且适合外星文明居住的行星。} \label{fig_Eaan_3}
\end{figure}
\begin{figure}[ht]
\centering
\includegraphics[width=6cm]{./figures/4f1f34a4309a75e0.png}
\caption{地球纳米布沙漠中的沙丘（上），与土卫六上贝莱特地区的沙丘对比} \label{fig_Eaan_4}
\end{figure}
在 1858 年至 1920 年 之间，包括一些科学家在内的许多人认为火星与地球非常相似，只是更加干燥，拥有厚重的大气层、相近的自转轴倾角、轨道和季节，并且存在一个建造了巨大火星运河的火星文明。这些理论由乔瓦尼·夏帕雷利（Giovanni Schiaparelli）、珀西瓦尔·洛厄尔（Percival Lowell） 等人提出。因此，在文学作品中，火星常被描绘成类似于地球沙漠的红色星球。然而，“水手号”（Mariner，1965 年）和“维京号”（Viking，1975–1980 年）空间探测器获得的图像和数据揭示，火星实际上是一个贫瘠、布满撞击坑的世界$^{[4][5][6][7][8][9]}$。尽管如此，随着研究的不断深入，一些将火星与地球类比的观点仍然存在。例如，“火星海洋假说”起源于维京号任务，并在 20 世纪 80 年代得到普及$^{[10]}$。随着过去存在液态水的可能性被提出，生命或许曾在火星上起源，火星也因此再次被视为一个更接近地球的世界。

同样地，直到 20 世纪 60 年代，包括一些科学家在内的许多人仍认为金星是一个更温暖版本的地球，拥有厚重的大气层，可能是炎热而多尘，或是潮湿并覆盖着水云和海洋$^{[11]}$。在科幻作品中，金星常被描绘成与地球相似的世界，许多人还曾推测存在金星文明。然而，这些看法在 20 世纪 60 年代被推翻：最早的空间探测器获取的精确科学数据表明，金星是一个极端炽热的星球，其表面温度约为 462 °C（864 °F）$^{[12]}$，并被酸性大气所包围，表面气压高达 9.2 MPa（1330 psi）$^{[12]}$。

自 2004 年起，“卡西尼–惠更斯号”（Cassini–Huygens）任务逐步揭示，土星的卫星泰坦是宜居带之外最类似地球的天体之一。尽管其化学组成与地球截然不同，但 2007 年对泰坦湖泊、河流以及流水地貌过程的确认，使其与地球的类比研究取得重大进展$^{[13][14]}$。后续观测（包括气象现象）进一步加深了人们对类地行星上可能发生的地质过程的理解$^{[15]}$。

开普勒空间望远镜自 2011 年起开始观测宜居带内潜在类地行星的凌日现象$^{[16][17]}$。尽管该技术为行星的发现和确认提供了更有效的手段，但仍无法最终判定这些候选行星在多大程度上类似地球$^{[18]}$。到 2013 年，多颗半径小于1.5 个地球半径、位于恒星宜居带内的开普勒候选行星被确认。直到 2015 年，首颗围绕类太阳恒星运行、尺寸接近地球的候选行星 Kepler-452b才被正式公布$^{[19][20]}$。

2023 年 1 月 11 日，NASA 科学家报告探测到 LHS 475 b——一颗类地系外行星，也是詹姆斯·韦布空间望远镜发现的第一颗系外行星$^{[21]}$。
\subsection{属性与判据}
寻找“地球类比体”的概率主要取决于被认为应当相似的属性，而这些属性在不同研究中差异很大。通常认为，这样的天体应是一颗类地行星，已有多项科学研究以此为目标展开探索。常被隐含但不限于的判据包括：行星大小、表面重力、恒星大小与类型（如类太阳恒星）、轨道距离及其稳定性、自转轴倾角与自转周期、相似的地理环境、海洋、大气与气候条件、强磁层，甚至是否存在类似地球的复杂生命。如果存在复杂生命，陆地上可能覆盖着大片森林；如果存在智慧生命，部分陆地甚至可能被城市所覆盖。然而，其中一些被假定的条件，考虑到地球自身的演化历史，可能并不常见。例如，地球的大气并非始终富含氧气，而富氧大气本身正是光合作用生命出现后的生物标志。此外，月球的形成、存在及其对这些特征的影响（如潮汐力）也可能使寻找地球类比体变得更加困难。

确定地球类比体的过程，往往需要在多个不确定性量化层面之间进行权衡。正如人类学家文森特·伊亚伦蒂（Vincent Ialenti）在其关于类比推理认识论的研究中所指出的那样$^{[22]}$，一些行星科学家“更愿意跨越时间与空间，凭借某种信念将两个迥异的对象联系起来”，而另一些人则对此持更为谨慎的态度$^{[23]}$。
\subsubsection{大小}
\begin{figure}[ht]
\centering
\includegraphics[width=8cm]{./figures/944eb86e695ec56c.png}
\caption{大小对比：Kepler-20e$^{[24]}$ 和 Kepler-20f$^{[25]}$ 与金星和地球的比较} \label{fig_Eaan_5}
\end{figure}
大小通常被认为是一个重要因素，因为与地球大小相近的行星更有可能在本质上属于类地行星，并且能够保留类似地球的大气层。$^{[26]}$

该列表包括质量位于 $0.8$–$1.9$ 个地球质量范围内的行星，低于这一范围的通常被归类为次地球型行星（sub-Earth），而高于这一范围的则被归类为超级地球（super-Earth）。此外，仅纳入半径已知落在 $0.5$–$2.0$ 个地球半径范围内（即地球半径的一半到两倍）的行星。

根据尺寸判据，按已知半径或质量来看，最接近地球的行星质量天体包括：
\begin{figure}[ht]
\centering
\includegraphics[width=14.25cm]{./figures/40bad47d87005cf0.png}
\caption{} \label{fig_Eaan_6}
\end{figure}
翻译“这一对比表明，仅凭大小作为衡量标准是不够的，尤其是在宜居性方面。还必须考虑温度因素，因为金星以及2012年发现的半人马座αB行星、开普勒-20（2011年发现））、COROT-7（2009年发现）以及开普勒-42的三颗行星（均于2011年发现）都极为炽热，而火星、木卫三和土卫六则是寒冷的世界，这也导致了它们表面和大气条件的极大差异。这些天体的质量为太阳系的卫星数量仅占地球卫星数量的一小部分，而系外行星的质量则很难准确测量。然而，发现类地行星——其大小与地球相当——具有重要意义，因为这可能表明类地行星的出现频率和分布情况。
\subsubsection{类地}
\begin{figure}[ht]
\centering
\includegraphics[width=6cm]{./figures/5c2efcf45e09b95e.png}
\caption{像土星卫星泰坦这样的表面（由“惠更斯”号探测器拍摄）在外观上与地球的洪泛平原具有一定的相似性。} \label{fig_Eaan_7}
\end{figure}
另一个常被提及的标准是：类地类比体必须是类地行星，也就是说，它应当具有相似的表面地质特征——由相似表面物质构成的行星表面。已知最接近的例子是火星和泰坦（土星的卫星）。尽管它们在地貌类型和表面组成上与地球存在一定相似性，但在温度和冰的含量等方面也存在显著差异。

地球上的许多表面物质和地貌，都是水参与作用的结果（例如黏土和沉积岩），或是生命活动的副产物（例如石灰岩或煤），也包括与大气的相互作用、火山活动，甚至人类活动的影响。因此，一个真正的类地类比体，可能需要经历类似的形成过程：拥有大气层，存在火山作用与表面的相互影响，曾经或正在存在液态水，以及生命形式。
\subsubsection{温度}
决定行星温度的因素有很多，因此在大气条件未知的情况下，也存在多种可用于与地球进行比较的温度指标$^{[citation\ needed]}$。对于没有大气层的行星，通常使用平衡温度；对于有大气层的行星，则假定存在温室效应；此外，还可以使用表面温度。这些温度指标都受到气候的影响，而气候又取决于行星的轨道与自转状态（或潮汐锁定），每一项都会引入额外的变量。

下面给出了与地球温度最为接近的已确认行星的比较。
\begin{figure}[ht]
\centering
\includegraphics[width=14.25cm]{./figures/e654f48183eac389.png}
\caption{} \label{fig_Eaan_8}
\end{figure}
\subsubsection{太阳类似体}
理想的、能够孕育生命的类地类比体的另一个标准是：它应当环绕一颗太阳类似体运行，即一颗与太阳十分相似的恒星。然而，这一标准并不一定完全成立，因为多种不同类型的恒星都可能提供适合生命存在的局部环境。例如，在银河系中，大多数恒星都比太阳更小、更暗。其中一颗这样的恒星是 TRAPPIST-1，距离地球约 12 秒差距（39 光年），其体积大约只有太阳的十分之一，亮度仅为太阳的约二千分之一，但它的宜居带内至少拥有六颗类地行星。尽管这些条件对已知生命而言似乎并不理想，但预计 TRAPPIST-1 仍将继续进行核燃烧约 12 万亿年（相比之下，太阳的剩余寿命约为 50 亿年），这为生命通过自然发生过程（无生源发生）起源提供了充足的时间$^{[36]}$。作为对比，地球上的生命在大约 10 亿年内便已演化出现$^{[citation\ needed]}$。
\subsubsection{地表水与水文循环}
\begin{figure}[ht]
\centering
\includegraphics[width=8cm]{./figures/4cc89536ed0e1272.png}
\caption{水覆盖了地球表面的约 70\%，并且是所有已知生命所必需的。} \label{fig_Eaan_9}
\end{figure}
\begin{figure}[ht]
\centering
\includegraphics[width=8cm]{./figures/aea374aae943db10.png}
\caption{位于类太阳恒星宜居带内的开普勒-22b，可能是迄今为止发现的、最有可能在其表面存在外星液态水的系外行星候选者之一，但它的体积明显大于地球，其真实组成仍然未知。} \label{fig_Eaan_10}
\end{figure}
宜居带（或称“液态水带”）这一概念，用于界定水能够在行星表面存在的区域，它基于地球与太阳的物理性质。在这一模型下，地球大致位于该区域的中心位置，即所谓的“金发姑娘位置”。地球是目前唯一被确认拥有大规模地表液态水的行星。金星位于宜居带偏热的一侧，而火星则位于偏冷的一侧。两者目前都未发现持续存在的地表水，不过有证据表明火星在远古时期曾经存在过液态水$^{[37][38][39]}$，而金星也被推测可能曾有类似情况$^{[11]}$。因此，位于“金发姑娘位置”、且拥有较厚大气层的系外行星（或卫星），可能像地球一样拥有海洋和水云。除地表水之外，一个真正的“地球类比体”还需要同时具备海洋或湖泊，以及未被水覆盖的陆地区域。

有观点认为，一个真正的地球类比体不仅需要在行星系统中处于相似的位置，还必须绕一颗类太阳恒星运行，并且其轨道接近圆形，从而像地球一样能够长期保持宜居状态$^{[\mathrm{citation\ needed}]}$。
\subsection{系外地球类比体}
“平庸原理”认为，偶然事件有可能使类似地球的行星在宇宙其他地方形成，并进而孕育出复杂的多细胞生命。与之相对，“稀有地球假说”则主张，如果采用最严格的判据，这样的行星即便存在，也可能遥远到人类永远无法发现。

由于太阳系中并不存在真正的地球类比体，相关探索已经扩展到系外行星领域。天体生物学家认为，地球类比体最有可能出现在恒星的宜居带内，在那里液态水可以存在，从而提供支持生命的基本条件。一些天体生物学家（如 Dirk Schulze-Makuch）还估计，质量足够大的天然卫星也可能形成类似地球的宜居卫星。
\subsubsection{历史}
\textbf{估计频率}
\begin{figure}[ht]
\centering
\includegraphics[width=8cm]{./figures/68c9d879b700c82e.png}
\caption{类地行星的艺术概念图 $^{[40]}$} \label{fig_Eaan_11}
\end{figure}
类地行星在银河系乃至整个宇宙中的出现频率目前仍不清楚，其估计范围从极端的“稀有地球假说”——仅有一个（即地球）——到数量几乎不可计数不等。

目前已有多项科学研究（包括开普勒任务）正利用凌日行星的真实观测数据来细化这一估计。亚利桑那大学天文学家 Michael Meyer 于 2008 年开展的一项研究，分析了新近形成的类太阳恒星附近的宇宙尘埃，结果表明，大约有 20\%–60\% 的太阳类恒星显示出形成类岩质行星的证据，其过程与地球形成的过程并无太大差异 $^{[41]}$。Meyer 团队在恒星周围发现了宇宙尘埃盘，并将其视为岩质行星形成过程的副产物。

2009 年，卡内基科学研究所的 Alan Boss 推测，仅在银河系中就可能存在多达 1000 亿颗类地行星 $^{[42]}$。

2011 年，美国国家航空航天局喷气推进实验室（JPL）基于开普勒任务的观测结果提出，约有 1.4\%–2.7\% 的类太阳恒星在其宜居带内拥有地球大小的行星。这意味着，仅银河系中就可能存在多达 20 亿颗地球大小的行星；如果假设其他星系中此类行星的数量与银河系相近，那么在可观测宇宙中的约 500 亿个星系内，类地行星的总数可能高达 10²⁰ 量级（即一百垓）$^{[43]}$。这一估计相当于在地球表面的每平方厘米上，平均对应约 20 个“地球类比体” $^{[44]}$。

2013 年，哈佛—史密森天体物理中心利用更多开普勒数据进行统计分析后指出，银河系中至少存在 170 亿颗地球大小的行星 $^{[45]}$。不过，这一结果并未涉及这些行星是否位于宜居带内。

2019 年的一项研究进一步认为，大约每 6 颗类太阳恒星中，就可能有 1 颗拥有地球大小的行星 $^{[46]}$。
\subsection{地球化改造}
\begin{figure}[ht]
\centering
\includegraphics[width=6cm]{./figures/5aacee2252da45f4.png}
\caption{经过地球化改造后的金星艺术想象图，被视为一种潜在的地球类比体} \label{fig_Eaan_12}
\end{figure}
地球化改造（Terraforming，字面意为“塑造地球”）是一个假想过程，指有意地改造某一行星、卫星或其他天体的**大气、温度、地表形态或生态系统**，使其接近地球的条件，从而适合人类居住。

由于距离较近且体积相似，火星被认为是最有可能进行地球化改造的候选天体，其次是金星，但可能性较低 $^{[47][48][49][50][51][52][53][54]}$。
\subsection{另见}
\begin{itemize}
\item 行星宜居性
\item 天然卫星的宜居性
\item 超宜居行星
\end{itemize}
\subsection{参考文献}
\begin{enumerate}
\item 奥布赖，丹尼斯（2013年11月4日）。“类地遥远行星遍布银河系”。《纽约时报》。检索日期：2013年11月5日
\item 佩蒂古拉，埃里克·A.；霍华德，安德鲁·W.；马西，杰弗里·W.（2013年11月1日）。“围绕类太阳恒星运行的地球大小行星的普遍性”. 美国国家科学院院刊. 110 (48): 19273–19278. arXiv:1311.6806. 文献标识码:2013PNAS..11019273P. DOI:10.1073/pnas.1319909110. PMC 3845182. PMID 24191033.
\item 舒尔策-马库赫，迪尔克；海勒，雷内；吉南，爱德华（2020年9月18日）。“寻找比地球更宜居的行星：超宜居世界的热门候选者”. 天体生物学. 20 (12): 1394–1404. 文献标识码:2020AsBio..20.1394S. DOI:10.1089/ast.2019.2161. PMC 7757576. PMID 32955925.
\item 奥加拉赫，J.J.；辛普森，J.A.（1965年9月10日）。“水手4号探测火星上被捕获电子与磁矩的搜索”。科学。新系列。149 (3689): 1233–1239. 天文学目录代码:1965Sci...149.1233O. doi:10.1126/science.149.3689.1233. PMID 17747452. S2CID 21249845.
\item 史密斯，爱德华·J.；戴维斯二世，莱弗雷特；科尔曼二世，保罗·J.；琼斯，道格拉斯·E.（1965年9月10日）。“火星附近磁场测量”。科学。新系列。149 (3689): 1241–1242. 文献标识码:1965Sci...149.1241S. doi:10.1126/science.149.3689.1241. PMID 17747454. S2CID 43466009.
\item 莱顿，罗伯特·B.；穆雷，布鲁斯·C.；夏普，罗伯特·P.；艾伦，J·丹顿；斯隆，理查德·K.（1965年8月6日）。“水手4号拍摄的火星照片：初步结果”。科学。新系列。149 (3684): 627–630. 文献标识码:1965Sci...149..627L. doi:10.1126/science.149.3684.627. PMID 17747569. S2CID 43407530.
\item 克利奥雷，阿尔维达斯；凯恩，丹·L.；列维，杰拉尔德·S.；埃什勒曼，冯·R.；菲耶尔博，冈纳尔；德雷克，弗兰克·D.（1965年9月10日）。“掩星实验：首次直接测量火星大气层与电离层的结果”。科学。新系列。149 (3689): 1243–1248. 文献标识码:1965Sci...149.1243K. DOI:10.1126/science.149.3689.1243. PMID 17747455. S2CID 34369864.
\item 索尔兹伯里，弗兰克·B.（1962年4月6日）。“火星生物学”。科学。新系列。136 (3510): 17–26. 文献标识码:1962Sci...136...17S. DOI:10.1126/science.136.3510.17. PMID 17779780. S2CID 39512870.
\item 基尔斯顿，史蒂文·D.；德拉蒙德，罗伯特·R.；萨根，卡尔（1966）。“以千米分辨率搜寻地球上的生命”。国际行星科学杂志. 5 (1–6): 79–98. 天文学数据库标识符:1966Icar....5...79K. DOI:10.1016/0019-1035(66)90010-8.
\item 美国国家航空航天局——火星海洋假说 已归档2012-02-20 于互联网档案馆
\item 桥本，G. L.；罗思-塞罗特，M.；杉田，S.；吉尔莫尔，M. S.；坎普，L. W.；卡尔森，R. W.；贝恩斯，K. H.（2008）。“伽利略近红外测绘光谱仪数据揭示的金星长英质高地地壳”. 地球物理研究期刊：行星. 113（E9）：E00B24。文献编码:2008JGRE..113.0B24H. DOI:10.1029/2008JE003134.
\item “金星概况”。于2016年3月8日从原始网页存档。检索日期：2016年3月10日。
\item “卡西尼号揭示土卫六的扎纳杜地区为类地地貌”。科学日报。2006年7月23日。检索日期：2007年8月27日。
\item “亲眼目睹、亲手触摸并亲鼻嗅闻土卫六这颗与地球极为相似的奇异世界”。欧空局新闻，欧洲航天局。2005年1月21日。检索日期：2005年3月28日。
\item “卡西尼-惠更斯：新闻”。Saturn.jpl.nasa.gov。于2008年5月8日从原始网页存档。检索日期：2011年8月20日。
\item “NASA开普勒望远镜证实首颗位于类太阳恒星宜居带的行星”。NASA新闻稿，2011年12月5日。于2017年5月16日存档自原文。检索日期：2011年12月6日
\item 霍威尔，伊丽莎白（2017年11月15日）。“开普勒-22b：宜居带系外行星的事实”。Space.com。于2019年8月22日从原文存档。检索日期：2019年2月10日
\item 佩蒂古拉，E. A.；霍华德，A. W.；马西，G. W.（2013）。“围绕类太阳恒星运行的地球大小行星的普遍性”. 美国国家科学院院刊. 110 (48): 19273–19278. arXiv:1311.6806. 天文学文献标识码:2013PNAS..11019273P. DOI:10.1073/pnas.1319909110. ISSN 0027-8424. PMC 3845182. PMID 24191033.
\item 詹金斯，乔恩·M.；特维肯，约瑟夫·D.；巴塔利亚，娜塔莉·M.等（2015年7月23日）。“开普勒-452b的发现与验证：一颗位于G2型恒星宜居带内的1.6倍地球半径超地球系外行星”。天文学杂志. 150 (2): 56. arXiv:1507.06723. 文献标识码:2015AJ....150...56J. DOI:10.1088/0004-6256/150/2/56. ISSN 1538-3881. S2CID 26447864.
\item “美国国家航空航天局望远镜在恒星的‘宜居带’发现一颗类地行星”。BNO新闻，2015年7月23日。于2016年3月4日从原文存档。检索日期：2015年7月23日
\item 周丹妮（2023年1月11日）“詹姆斯·韦伯望远镜发现首颗系外行星——据约翰霍普金斯大学应用物理实验室天文学家领导的研究团队称，这颗行星的大小几乎与地球相当”。NBC新闻。已获取2023年1月12日
\item 深时反思——一个星球。麻省理工学院出版社，2020年9月22日。ISBN978-0-262-53926-5。检索日期：2023年1月14日。
\item 莫根，菲利普（2023年1月12日）。“在地球上寻找外星世界”：引用期刊需（帮助)
\item 美国国家航空航天局工作人员（2011年12月20日）。“开普勒：寻找宜居行星——开普勒-20e”。美国国家航空航天局。于2012年3月10日从原始网页存档。检索日期：2011年12月23日
\item 美国国家航空航天局工作人员（2011年12月20日）。“开普勒：寻找宜居行星——开普勒-20f”。美国国家航空航天局。于2012年3月10日从原始网页存档。检索日期：2011年12月23日
\item 埃尔卡耶夫，N.V.；拉默，H.；埃尔金斯-坦顿，L.T.；施托克尔，A.；奥德特，P.；马克，E.；多尔菲，E.A.；基斯利亚科娃，K.G.；库利科夫，Yu.N.；莱茨inger，M.；居德尔，M.（2014）。“火星原始大气层的逃逸与初始水储量”。行星与空间科学98：106–119。arXiv1308.0190文献标识码2014P&SS...98..106EDOI10.1016/j.pss.2013.09.008ISSN0032-0633PMCID4375622PMID25843981
\item 托雷斯，吉列尔莫；弗雷桑，弗朗索瓦（2011）。“将开普勒凌星光变曲线建模为假阳性：排除开普勒-9的混叠情形，并验证开普勒-9d——一颗位于多行星系统中的超地球大小行星”。天体物理学杂志. 727 (24): 24. arXiv:1008.4393. Bibcode:2011ApJ...727...24T. doi:10.1088/0004-637X/727/1/24. S2CID 6358297.
\item 约翰逊，米歇尔；哈林顿，J.D.（2014年4月17日）。“NASA的开普勒望远镜首次发现一颗位于另一颗恒星‘宜居带’内的类地行星”。NASA。检索日期：2014年4月17日
\item 布鲁格，B.；等（2016年11月10日）。“比邻星b的可能内部结构与组成”. 天体物理学报快报. 831（2）：L16。arXiv:1609.09757. 文献编码:2016ApJ...831L..16B. doi:10.3847/2041-8205/831/2/L16.
\item 约翰逊，米歇尔（2011年12月20日）。“美国国家航空航天局发现首个位于太阳系外、与地球大小相当的行星”。美国国家航空航天局。于2019年5月4日存档自原文。检索日期：2011年12月20日
\item 汉德，埃里克（2011年12月20日）。“开普勒发现首批类地系外行星”。自然. doi:10.1038/nature.2011.9688. S2CID 122575277.
\item 致恒星主序
\item “美国国家航空航天局，火星：事实与数据”。于2004年1月23日从原始网页存档。检索日期：2010年1月28日。
\item 马拉马，A.；王，D.；霍华德，R.A.（2006）。“金星相函数与HSO的前向散射”。国际行星科学杂志。182（1）：10–22。文献标识码：2006Icar..182...10Mdoi10.1016/j.icarus.2005.12.014
\item 马拉马，A.（2007）。“火星的亮度与反照率”。国际行星科学杂志. 192 (2): 404–416. 天文学文献标识符:2007Icar..192..404M. DOI:10.1016/j.icarus.2007.07.011.
\item 斯内伦，伊格纳斯·A·G.（2017年2月）。“地球的七大姐妹”. 自然. 542 (7642): 421–422. doi:10.1038/542421a. hdl:1887/75076. PMID 28230129.
\item 卡布罗尔，N. 和 E. 格林（编）。2010年。《火星上的湖泊》。爱思唯尔。纽约
\item 克利福德，S. M.；帕克，T. J.（2001）。“火星水圈的演化：对原始海洋命运及北区平原现状的启示”。伊卡洛斯.154(1)：40–79。Bibcode：2001Icar..154...40Cdoi10.1006/icar.2001.6671
\item 比利亚努埃瓦，G.；芒马，M.；诺瓦克，R.；考夫尔，H.；哈托赫，P.；昂克雷纳兹，T.；德仓永，A.；卡亚特，A.；史密斯，M.（2015）。“火星大气中强烈的水同位素异常：探究当前与古代储库”. 科学. 348 (6231): 218–221. 文献编码:2015Sci...348..218V. DOI:10.1126/science.aaa3630. PMID 25745065. S2CID 206633960.
\item "Artist's Concept of Earth-Like Planets in the Future Universe". ESA/Hubble. Retrieved 22 October 2015.
\item Briggs, Helen (17 February 2008). "Planet-hunters set for big bounty". BBC News.
\item 帕沃夫斯基，A.（2009年2月25日）。“银河系可能充满‘地球’般的行星，孕育外星生命”. CNN.
\item 崔，查尔斯·Q.（2011年3月21日）。“外星宜居行星新估算：仅我们银河系就有20亿颗”。Space.com。检索日期：2011年4月24日
\item “Wolfram|Alpha：让世界知识可计算”.www.wolframalpha.com。检索日期：2018年3月19日
\item 银河系中存在170亿颗类地大小的外星行星 已归档2014年10月6日，存于互联网档案馆SPACE.com 2013年1月7日
\item 徐丹莉·C.；福特埃里克·B.；拉戈齐内达林；阿什比凯尔（2019年8月14日）。“围绕FGK型恒星运行的行星出现率：结合开普勒DR25、盖亚DR2与贝叶斯推断”. 天文学杂志. 158 (3): 109. arXiv:1902.01417. 文献标识码:2019AJ....158..109H. DOI:10.3847/1538-3881/ab31ab. ISSN 1538-3881. S2CID 119466482.
\item 罗伯特·M·祖布林（先驱航天公司），克里斯托弗·P·麦凯。美国国家航空航天局艾姆斯研究中心（约1993年）。“火星地球化的技术需求”.
\item 马特·康威（2007年2月27日）。“我们终于到了：火星 terraforming”。Aboutmyplanet.com。于2011年7月23日从原文存档。检索日期：2011年8月20日。”
\item 彼得·阿伦斯。“世界的地球化改造”（PDF）。Nexial Quest。于2019年6月9日从原文存档。检索日期：2007年10月18日。
\item 萨根，卡尔（1961）。“金星行星”. 科学. 133 (3456): 849–58. 文献标识码:1961Sci...133..849S. DOI:10.1126/science.133.3456.849. PMID 17789744.
\item 祖布林，罗伯特（1999）。迈向太空：打造一个航天文明。企鹅出版社。ISBN 9781585420360.
\item 福格，马丁·J.（1995）。地球化：行星环境工程。SAE国际，宾夕法尼亚州沃伦代尔。ISBN 1-56091-609-5.
\item 伯奇，保罗（1991）。“快速改造金星” （PDF）. 英国行星际学会期刊. 44: 157. 文献标识码:1991JBIS...44..157B.
\item 兰迪斯，杰弗里·A.（2003年2月2日至6日）。“金星的殖民化”. 人类太空探索会议、空间技术与应用国际论坛，新墨西哥州阿尔伯克基. 654: 1193. 文献标识码:2003AIPC..654.1193L. DOI:10.1063/1.1541418.
\end{enumerate}

